% ============================================================
\chapter{Th\'eor\`eme de Markov-Kakutani}
\label{ch:markov_kakutani}
% ============================================================

En 1936, Andrei Markov (fils du c\'el\`ebre probabiliste) et, ind\'ependamment, Shizuo Kakutani d\'emontrent un r\'esultat remarquable~: toute famille d'applications affines continues qui commutent, agissant sur un convexe compact d'un espace localement convexe, poss\`ede un point fixe commun. Ce th\'eor\`eme, dont l'\'enonc\'e semble purement g\'eom\'etrique, a des ramifications profondes en analyse harmonique abstraite~: il est \`a la base de l'existence de \emph{moyennes invariantes} sur les groupes ab\'eliens et joue un r\^ole fondateur dans la th\'eorie des \emph{groupes moyennables}, d\'evelopp\'ee par John von Neumann et poursuivie par Mahlon Day et Erling F\o lner.

\begin{intuition}
Le th\'eor\`eme de Markov-Kakutani garantit l'existence d'un point fixe commun pour une
famille d'applications affines continues qui commutent, agissant sur un convexe compact
d'un espace vectoriel topologique localement convexe. Ce r\'esultat, simple dans son
\'enonc\'e, a des applications profondes~: il est \`a la base de la th\'eorie des moyennes
invariantes et joue un r\^ole central dans l'\'etude des groupes moyennables.
\end{intuition}

% ------------------------------------------------------------
\section{\'Enonc\'e et preuve du th\'eor\`eme}
% ------------------------------------------------------------

\begin{definition}[Application affine]
\index{application affine}
Soit $E$ un espace vectoriel et $C \subseteq E$ un convexe. Une application $T\colon C \to C$
est \emph{affine} si pour tous $x, y \in C$ et tout $\lambda \in [0,1]$~:
\[
  T(\lambda x + (1-\lambda)y) = \lambda T(x) + (1-\lambda) T(y).
\]
\end{definition}

\begin{theoreme}[Markov-Kakutani]
\label{thm:markov-kakutani}
\index{th\'eor\`eme de Markov-Kakutani}
Soit $E$ un espace vectoriel topologique localement convexe et $C \subseteq E$ un convexe
compact non vide. Soit $\mathcal{F}$ une famille d'applications affines continues
$T\colon C \to C$ qui commutent deux \`a deux ($T_1 \circ T_2 = T_2 \circ T_1$ pour tous
$T_1, T_2 \in \mathcal{F}$). Alors il existe $x^* \in C$ tel que $T(x^*) = x^*$ pour
tout $T \in \mathcal{F}$.
\end{theoreme}

\begin{proof}
\emph{\'Etape 1~: cas d'une seule application.}
Soit $T\colon C \to C$ affine continue. Pour tout $x \in C$ et $n \geq 1$, posons
\[
  \sigma_n(x) = \frac{1}{n} \sum_{k=0}^{n-1} T^k(x).
\]
Comme $C$ est convexe et $T(C) \subseteq C$, on a $\sigma_n(x) \in C$. D\'efinissons
$C_n = \sigma_n(C) = \{\sigma_n(x) : x \in C\}$. Alors $C_n$ est compact (image continue
d'un compact) et non vide.

Montrons que $(C_n)_{n \geq 1}$ a la propri\'et\'e d'intersection finie, ce qui implique
$\bigcap_n \overline{C_n} \neq \emptyset$ par compacit\'e de $C$.

Soit $y \in \bigcap_n \overline{C_n}$. Pour tout $\varepsilon$-voisinage $V$ de $0$, il
existe $x_n \in C$ tel que $\sigma_n(x_n) - y \in V$. On v\'erifie que
$T(\sigma_n(x_n)) - \sigma_n(x_n) = \frac{1}{n}(T^n(x_n) - x_n) \to 0$ car $C$ est
born\'e. Ainsi $T(y) = y$.

\emph{\'Etape 2~: famille commutante.}
Pour chaque $T \in \mathcal{F}$, d\'efinissons $\mathrm{Fix}(T) = \{x \in C : T(x) = x\}$.
Par l'\'etape~1, $\mathrm{Fix}(T)$ est non vide. De plus~:
\begin{itemize}
  \item $\mathrm{Fix}(T)$ est ferm\'e (pr\'eimage du ferm\'e diagonal par $(x, T(x))$)
    et contenu dans le compact $C$, donc compact.
  \item $\mathrm{Fix}(T)$ est convexe car $T$ est affine~: si $T(x) = x$ et $T(y) = y$,
    alors $T(\lambda x + (1-\lambda)y) = \lambda x + (1-\lambda)y$.
\end{itemize}
Si $T_1, T_2$ commutent et $x \in \mathrm{Fix}(T_1)$, alors
$T_1(T_2(x)) = T_2(T_1(x)) = T_2(x)$, donc $T_2(\mathrm{Fix}(T_1)) \subseteq \mathrm{Fix}(T_1)$.
Ainsi $T_2$ restreinte \`a $\mathrm{Fix}(T_1)$ est une application affine continue d'un
convexe compact dans lui-m\^eme, et par l'\'etape~1, elle admet un point fixe dans
$\mathrm{Fix}(T_1)$.

Par une r\'ecurrence transfinie (ou par la propri\'et\'e d'intersection finie),
$\bigcap_{T \in \mathcal{F}} \mathrm{Fix}(T) \neq \emptyset$.
\end{proof}

% ------------------------------------------------------------
\section{Groupes moyennables}
% ------------------------------------------------------------

\begin{definition}[Moyenne invariante]
\index{moyenne invariante}
Soit $G$ un groupe discret. Une \emph{moyenne invariante \`a gauche} sur $G$ est une forme
lin\'eaire $m\colon \ell^\infty(G) \to \R$ telle que~:
\begin{enumerate}
  \item $m(\ind_G) = 1$ (normalisation)~;
  \item $m(f) \geq 0$ si $f \geq 0$ (positivit\'e)~;
  \item $m(L_g f) = m(f)$ pour tout $g \in G$ (invariance), o\`u $(L_g f)(x) = f(g^{-1}x)$.
\end{enumerate}
\end{definition}

\begin{definition}[Groupe moyennable]
\index{groupe moyennable}
Un groupe discret $G$ est \emph{moyennable} s'il admet une moyenne invariante \`a gauche.
\end{definition}

\begin{theoreme}[Les groupes ab\'eliens sont moyennables]
\label{thm:abelien-moyennable}
Tout groupe ab\'elien discret est moyennable.
\end{theoreme}

\begin{proof}
Soit $G$ un groupe ab\'elien. L'ensemble des moyennes
\[
  \mathcal{M} = \{m \in (\ell^\infty(G))^* : m \geq 0, \; m(\ind_G) = 1\}
\]
est un convexe compact non vide pour la topologie faible-$*$ (par le th\'eor\`eme de
Banach-Alaoglu).

Pour chaque $g \in G$, l'application $T_g\colon \mathcal{M} \to \mathcal{M}$ d\'efinie par
$(T_g m)(f) = m(L_{g^{-1}} f)$ est affine et continue (pour la topologie faible-$*$).
Elle envoie $\mathcal{M}$ dans $\mathcal{M}$ car $L_{g^{-1}}$ pr\'eserve la positivit\'e
et $L_{g^{-1}} \ind_G = \ind_G$.

Comme $G$ est ab\'elien, $L_g \circ L_h = L_{gh} = L_{hg} = L_h \circ L_g$, donc les
$T_g$ commutent. Par le th\'eor\`eme de Markov-Kakutani, il existe
$m^* \in \bigcap_{g \in G} \mathrm{Fix}(T_g)$, c'est-\`a-dire une moyenne invariante.
\end{proof}

\begin{exemple}[Moyenne invariante sur $\Z$]
La moyenne de Banach sur $\Z$, construite ci-dessus, n'est pas explicite. N\'eanmoins, on
peut montrer que pour toute suite born\'ee $(a_n)_{n \in \Z}$, si la limite de Ces\`aro
$\lim_{N \to \infty} \frac{1}{2N+1} \sum_{n=-N}^{N} a_n$ existe, alors elle co\"incide
avec $m^*((a_n))$.
\end{exemple}

% ------------------------------------------------------------
\section{Exemples de groupes non moyennables}
% ------------------------------------------------------------

\begin{theoreme}[von Neumann]
Le groupe libre \`a deux g\'en\'erateurs $F_2$ n'est pas moyennable.
\end{theoreme}

\begin{proof}[Id\'ee de la preuve]
On utilise le paradoxe de Banach-Tarski~: si $F_2$ \'etait moyennable, on pourrait
construire une mesure finiment additive invariante sur $F_2$, ce qui contredit le fait que
$F_2$ admet une d\'ecomposition paradoxale (partition en quatre morceaux qui, par
translations, reconstituent deux copies de $F_2$).
\end{proof}

\begin{proposition}[Propri\'et\'es de la moyennabilit\'e]
\begin{enumerate}
  \item Les sous-groupes d'un groupe moyennable sont moyennables.
  \item Les quotients d'un groupe moyennable sont moyennables.
  \item Les extensions de groupes moyennables par des groupes moyennables sont moyennables.
  \item Tout groupe fini est moyennable (prendre la moyenne uniforme).
  \item Un groupe contenant $F_2$ comme sous-groupe n'est pas moyennable.
\end{enumerate}
\end{proposition}

% ------------------------------------------------------------
\section{Variantes et g\'en\'eralisations}
% ------------------------------------------------------------

\begin{theoreme}[Ryll-Nardzewski]
\label{thm:ryll-nardzewski}
\index{th\'eor\`eme de Ryll-Nardzewski}
Soit $E$ un espace de Banach et $C \subseteq E$ un convexe faiblement compact non vide.
Soit $\mathcal{F}$ un semi-groupe d'isom\'etries affines de $C$. Alors
$\bigcap_{T \in \mathcal{F}} \mathrm{Fix}(T) \neq \emptyset$.
\end{theoreme}

\begin{remarque}
Le th\'eor\`eme de Ryll-Nardzewski g\'en\'eralise Markov-Kakutani en rempla\c{c}ant la
commutativit\'e par une condition de non-expansivit\'e et en affaiblissant la compacit\'e
en compacit\'e faible.
\end{remarque}

\begin{formules}[title=\textbf{Points fixes communs}]
\begin{center}
\renewcommand{\arraystretch}{1.3}
\begin{tabular}{|l|l|l|}
\hline
\textbf{Th\'eor\`eme} & \textbf{Famille} & \textbf{Espace} \\
\hline
Markov-Kakutani & Affines commutantes & Convexe compact, evtlc \\
Ryll-Nardzewski & Isom\'etries affines & Convexe faiblement compact, Banach \\
\hline
\end{tabular}
\end{center}
\end{formules}

% ------------------------------------------------------------
\section{Exercices}
% ------------------------------------------------------------

\begin{exercice}
Montrer que le th\'eor\`eme de Markov-Kakutani est faux sans l'hypoth\`ese de
commutativit\'e. \emph{Indication~: consid\'erer deux rotations du disque.}
\end{exercice}

\begin{exercice}
Appliquer le th\'eor\`eme de Markov-Kakutani pour montrer que tout groupe ab\'elien
compact $G$ admet une mesure de Haar (en prenant $C$ l'ensemble des mesures de probabilit\'e
sur $G$).
\end{exercice}

\begin{exercice}
Montrer que le groupe des isom\'etries du plan euclidien n'est pas moyennable.
\emph{Indication~: montrer qu'il contient $F_2$ comme sous-groupe.}
\end{exercice}

\begin{exercice}
Soit $T\colon C \to C$ une application affine continue sur un convexe compact $C$ d'un
espace de Banach. Montrer directement (sans Markov-Kakutani) que les moyennes de Ces\`aro
$\sigma_n(x) = \frac{1}{n}\sum_{k=0}^{n-1} T^k(x)$ convergent (au moins le long d'une
sous-suite) vers un point fixe de $T$.
\end{exercice}

\begin{exercice}
Soit $G$ un groupe moyennable discret agissant sur un espace compact $X$ par
hom\'eomorphismes. Montrer qu'il existe une mesure de probabilit\'e bor\'elienne sur $X$
invariante par l'action de $G$.
\end{exercice}

\begin{exercice}[Moyennabilit\'e et suites de F\o lner]
\index{suite de F\o lner}
Montrer qu'un groupe d\'enombrable $G$ est moyennable si et seulement s'il admet une
\emph{suite de F\o lner}~: une suite $(F_n)$ d'ensembles finis non vides de $G$ tels que
pour tout $g \in G$,
\[
  \frac{\abs{gF_n \triangle F_n}}{\abs{F_n}} \to 0 \quad \text{quand } n \to \infty.
\]
V\'erifier que $F_n = \{-n, \ldots, n\}$ est une suite de F\o lner pour $\Z$.
\end{exercice}
